\subsection{Theorem C: obstruction analysis and certificate method}\label{sec:methodC}
The analytic obstructions are proved by hand. Proposition~\ref{prop:Cupper}: the template declares $n=A$ losing, but for $k\ge A$ removing $A$ counters is a legal winning first move, a contradiction; hence any valid template needs $k<A$. Proposition~\ref{prop:Ccollapse}: a nonzero translation preserving the three-residue set must have order three, so the set is the orbit $\{0,p/3,2p/3\}$; this forces $A=p/3$, i.e.\ $k=A-1$, with displayed modulus $3A$ and least period $A$.
The twelve fixed-$k$ certificates use the Theorem~B window-and-bridge pipeline with $A=54$ via \texttt{methodology/theorem\_C/verify\_theorem\_c.py}, with the same Sprague--Grundy cross-checks (Table~\ref{tab:Ccert}). Template boundaries are proved as finite failures: the script searches the exact recurrence for the first heap size where the proposed $A$-template disagrees with the true starting outcome and confirms it against the naive solver (Table~\ref{tab:Cfail}).
